% Generated draft: requires mathematical and source review.
\documentclass[11pt]{article}
\usepackage[margin=25mm]{geometry}
\usepackage{fontspec}
\setmainfont{Latin Modern Roman}
\usepackage{amsmath,amssymb,mathtools,unicode-math}
\setmathfont{Latin Modern Math}
\usepackage{enumitem,needspace,xurl,hyperref}
\setlength{\parindent}{0pt}
\setlength{\parskip}{5pt}
\setlength{\emergencystretch}{3em}
\newcommand{\N}{\mathbb{N}}
\newcommand{\Z}{\mathbb{Z}}
\newcommand{\Q}{\mathbb{Q}}
\newcommand{\R}{\mathbb{R}}
\newcommand{\C}{\mathbb{C}}
\newcommand{\F}{\mathbb{F}}
\newcommand{\A}{\mathbb{A}}
\newcommand{\PP}{\mathbb{P}}
\newcommand{\E}{\mathbb{E}}
\newcommand{\eps}{\varepsilon}
\newcommand{\dd}{\,\mathrm{d}}
\begin{document}
\section*{UnsolvedMath ID: 20003189}
\textbf{Source problem code:} catalogue entry
\subsection{Definitions and mathematical statement}
Let $S$ be a closed oriented genus-$g\geq2$ surface and $H_n$ an $SL(n,\R)$ Hitchin component. The source writes $A=\Q(H_n)$ and a central extension $1\to K_2(A)\to E(n,A)\to SL(n,A)\to1$, and asks whether the logarithmic differential of the resulting surface class is a Weil--Petersson form. A precise corrected version first chooses a smooth irreducible algebraic moduli model $Y$ over a characteristic-zero field $k$, containing the intended Hitchin locus in its real points, and a generically descended tautological representation over $A=k(Y)$. An arbitrary constant $A$-valued representation would not represent the varying moduli symplectic form. A coarse character variety need not itself carry such a universal matrix representation without a framing, stack, or field extension.

For a specified central extension $1\to K_2(A)\to\widetilde G_A\to SL(n,A)\to1$, choose generators $a_i,b_i$ with $\prod_i[a_i,b_i]=1$ and lift their images under $\rho$. The product of lifted commutators gives
\[
q_\rho=\prod_{i=1}^g[\widetilde\rho(a_i),\widetilde\rho(b_i)]\in K_2(A),
\]
independent of the central choices of lifts, with orientation/sign conventions fixed. Over a field in the stable range $n\geq3$, a Steinberg extension realizes the usual $K_2$ kernel. Rank two needs an unstable extension or a specified compatible stabilization; the displayed exact sequence cannot silently be assumed to have the same stable kernel. The source's symbol $E(n,A)$ is treated as its extension notation, not automatically the customary elementary matrix group.

For Milnor symbols the logarithmic differential is
\[
d\log\{f,g\}=\frac{\dd f}{f}\wedge\frac{\dd g}{g}\in\Omega^2_{A/k}.
\]
Initially this is a rational algebraic form, not automatically a regular differential form on all of $Y$ or $H_n$. Trivial codimension-one tame symbols of the $K_2$ class supply the unramified regularity input in the smooth model. The aid reports that Fock--Goncharov's 2006 motivic surface class gives the generalized Weil--Petersson/Goldman form after matching the algebraic model, central extension, orientation, and invariant quadratic-form normalization. The three stated checks are generic tautological descent, codimension-one regularity, and normalization compatibility; their verification for a literal reading of the AIM notation is not supplied.

With trace pairings $B_n(U,V)=\operatorname{Tr}_n(UV)$, the Goldman form is
\[
\omega_n([\alpha],[\beta])=\langle B_n(\alpha\smile\beta),[S]\rangle
\]
on the smooth representation-moduli locus. For the principal Fuchsian embedding $F_n$ induced by the irreducible $n$-dimensional representation of $SL(2)$, the aid gives the exact test
\[
F_n^*\omega_n=\frac{n(n^2-1)}6\,\omega_2.
\]
Identifying $\omega_2$ with the conventionally normalized Weil--Petersson form requires its scalar normalization as well. Thus the literal question is not a normalization-free equality, and a higher-rank form cannot simply restrict with factor one under all trace conventions. The corrected motivic comparison is a supplied literature assertion with missing full theorem details, not an independently verified identity for every possible model or extension.
\subsection{Short English statement}
The supplied corrected answer identifies the surface $K_2$ class with the Goldman/generalized Weil--Petersson form only after specifying a tautological algebraic model, regularity, and matching extension normalization. Rank two needs an additional extension caveat. The principal Fuchsian restriction has trace-normalization factor $n(n^2-1)/6$, so an unqualified scalar-free equality would be misleading.
\subsection{Sources}
\begin{itemize}
\item Local registry: permanent catalogue entry, exact source code catalogue entry; \texttt{aim-topology-notes.json}, \texttt{notes[276]}; canonical tag ``question''.
\item AIM workshop \emph{Representations of surface groups}, Question 4.7 attributed to Olivier Guichard, \url{https://aimath.org/WWN/surfacegroups/surfacegroups.pdf}. It refers to an omitted preceding question. No complete original literature citation, proposer year, or remarks are supplied.
\item Statement status is \texttt{corrected\_verified}. Verification reports separating the central-extension display from ``SL(n)'' and removing the next heading ``5 Other questions''; the mathematical reconstruction follows that inherited check.
\item Supplied reference: Fock--Goncharov's 2006 motivic surface-class construction, without title, theorem number, or complete authoritative publication details.
\item Supplied research attempt 1, registry updated 14 August 2026: partial result, low novelty confidence, machine-generated and not peer reviewed. It supplies the descent/regularity/normalization audit and principal trace-index test.
\end{itemize}
\end{document}
